\begin{claim}{C.37}{Proved}{Uniform control on all two-dimensional operator closures}
Let $f_t$, $0\leq t\leq1$, be an ambient-null based loop in
$\operatorname{Homeo}^+(\R^2)$, let $G=\langle f_t:0\leq t\leq1\rangle$,
and let $H$ be any subgroup commuting with every $f_t$.
For an open $G$-orbit $O$, let $Q_O$ be the actual closed restricted
operator image of C.30. Put
\[
 U_2=\bigcup_{\dim Q_O=2}O,\qquad X_2=\overline{U_2}.
                                                        \tag{C.37.1}
\]
The contractions of C.36 extend to a jointly continuous based
contraction through homeomorphisms of $X_2$, with jointly continuous
inverses, commuting with $H$. This includes arbitrary accumulation
of the open orbits in (C.37.1).

The uniform estimate used in the proof is the following exact
word statement. In either simply connected two-dimensional Lie
group $\R^2$ or $\operatorname{Aff}^+(\R)$, every continuous path
$k:[0,t]\to N$ from $1$ to $g$ satisfies
\[
 \exp(r\log g)
   =k(b)k(a)^{-1}
 \quad\hbox{or}\quad
 \exp(r\log g)=g\,k(b)k(a)^{-1}
 \quad(0\leq r\leq1)                                   \tag{C.37.2}
\]
for suitable $a,b\in[0,t]$. The same assertion is used in the
universal covering group when the actual group is $\R\times S^1$.
Consequently every contraction value on an open orbit in $U_2$
is represented by a word of length at most three in the original
loop and its inverse, with equality of the specified based
universal-cover lifts.
\end{claim}

\begin{proof}
\emph{The three-word estimate.}
For a continuous path $\gamma:[0,1]\to\R^2$ with
$\gamma(0)=0$, $\gamma(1)=v$, one has
\[
 rv=\gamma(b)-\gamma(a)
 \quad\hbox{or}\quad
 rv=v+\gamma(b)-\gamma(a) \qquad(0\leq r\leq1).           \tag{C.37.3}
\]
Here is an exact measure-theoretic proof that applies to arbitrary
paths. For a polygonal path apply the signed-measure splitting
theorem of Grandoni--Ravi--Singh--Zenklusen
\cite[Theorem 4.7 and \S4.3.2]{C_GrandoniRaviSinghZenklusen2014}
to the two signed measures with densities $\gamma_1'$ and
$\gamma_2'$. A single circular interval has fraction $r$ of both
total measures. Its integral is either a path difference or the
total vector plus a path difference, giving (C.37.3).
Approximate an arbitrary continuous path uniformly by polygonal
paths with the same endpoints. Compactness of the two endpoints
and of the two alternatives passes the equation to the limit.
The planar case of the cited splitting proof uses at most two
cuts of an interval; the two equal-mass halves together have at
most three pieces, so one is a single interval. This is the
one-circular-interval count required here.

For $N=\R^2$, (C.37.2) is (C.37.3).
Write the affine group in the exact coordinates of C.36:
\[
 (a,b)(a',b')=(a+a',\,b+e^a b').
                                                        \tag{C.37.4}
\]
Every nonidentity $g=(a,b)$ lies on its unique exponential
one-parameter subgroup. If $a\ne0$, conjugation by
$(0,b/(e^a-1))$ sends $g$ to $(a,0)$.
The global coordinate map
\[
 (a',b')\longmapsto (a'/a,\,e^{-a'}b')
                                                        \tag{C.37.5}
\]
then sends left multiplication by $g^r$ to translation by
$(r,0)$, and sends $1,g$ to $(0,0),(1,0)$.
If $g=(0,b)$ with $b\ne0$, use $(a',b')\mapsto(b'/b,a')$
instead. Apply (C.37.3) in these coordinates.
A difference $(r,0)$ says $k(b)=g^r k(a)$; the other
alternative says $k(a)=g^{1-r}k(b)$. These are exactly
(C.37.2). The identity case uses the identity word.
This proves the estimate also for the simply connected
covering group of every group in C.36.

\emph{The compact family of actual word isotopies.}
Use $r=1-s$ for the contraction parameter, and set
\[
 R^O_{r,t}=p_O\exp\bigl(r\log\widetilde n_O(t)\bigr),
 \qquad n_O(t)=f_t|_O,                                  \tag{C.37.6}
\]
as in C.36. Consider the following compact family of actual
plane isotopies, with the compact-open topology on maps and
inverses:
\[
 W_{\epsilon,a,b,t}(u)
       = f_{ut}^{\,\epsilon} f_{ub}f_{ua}^{-1},
 \quad 0\leq a,b\leq t\leq1,\quad
       \epsilon\in\{0,1\},\quad 0\leq u\leq1 .
                                                        \tag{C.37.7}
\]
Compactness follows from continuity of the original loop and
of multiplication and inversion. Formula (C.37.2), applied to
the based lift of its prefix $[0,t]$, represents every
$R^O_{r,t}$ by the endpoint of one of (C.37.7).
More precisely, the based lift of that word isotopy in $N_O$
has endpoint
$\exp(r\log\widetilde n_O(t))$. Thus the representation retains
the lift, not only the projected homeomorphism.

\emph{The exact groups on nonopen limit orbits.}
Each $N_O$ in (C.37.6) is metabelian. Therefore the identity
\[
          [[g_1,g_2],[g_3,g_4]]=1
                                                        \tag{C.37.8}
\]
holds on $U_2$ for all $g_i\in G$, and hence on $X_2$ by
continuity. The set $X_2$ is $G$-invariant, so it is a union
of whole $G$-orbits. Its open $G$-orbits are precisely those
already in $U_2$: any other open orbit has an open
neighborhood disjoint from $U_2$.

Let $B\subset X_2$ be a nontrivial nonopen orbit, endowed
with the intrinsic line or circle topology of C.28.
The image $G_B$ in $\operatorname{Homeo}^+(B)$ is transitive,
metabelian, and path connected. For the last assertion,
every finite word has a path from the identity obtained from
the original paths and their inverses; C.28 gives intrinsic
compact-open continuity of all these finite-word paths.
We now prove the exact classification needed at $B$:
\[
 G_B\cong
 \begin{cases}
   (\R,+)\ \hbox{or }\operatorname{Aff}^+(\R),& B\text{ a line},\\
   S^1,& B\text{ a circle},
 \end{cases}                                            \tag{C.37.9}
\]
with the usual transitive actions after conjugacy.

First, a transitive abelian subgroup $A$ of
$\operatorname{Homeo}^+(B)$ has abelian closure $\overline A$,
and this closure acts freely: an element fixing one point
commutes with $A$, so fixes its whole transitive orbit.
The closure is Polish and transitive. The open orbit-map
theorem used in C.26 identifies it homeomorphically with $B$
by evaluation, since its point stabilizer is trivial.
In particular it is locally compact. Its faithful product
action on the three-manifold $B\times\R^2$ makes it a Lie
group by the same Hilbert--Smith input as C.33.
It is the one-dimensional connected group $\R$ or $S^1$,
and evaluation conjugates its action to translation.
Since $A$ was already transitive, injectivity of evaluation
also proves $A=\overline A$.

If $G_B$ is nonabelian, its derived subgroup $D$ is
nontrivial, abelian, normal, and path connected.
Path connectedness follows by taking commutator paths and
finite products of them. The common fixed set of $D$ is
$G_B$-invariant. Transitivity and nontriviality make it empty.
Every $D$-orbit is therefore a connected nontrivial interval
or circle. Such an orbit has no included endpoint, since
that endpoint would be preserved by its transitive
orientation-preserving action. All these orbits are open in
$B$; connectedness of $B$ forces $D$ to be transitive.
The preceding paragraph identifies $D$ with the full
translation group.

The orientation-preserving normalizer of translations on
the line consists exactly of $x\mapsto cx+d$, $c>0$.
Indeed normalization says $h(x+a)=h(x)+\chi(a)$;
continuity and additivity give $\chi(a)=ca$.
The multiplier image of the path-connected group $G_B$
is a connected subgroup of $\R_{>0}$. It is nontrivial in
the present case, so it is all $\R_{>0}$; because $G_B$
contains every translation, it is the full affine group.
On the circle, the orientation-preserving normalizer of
the full rotation group is just that rotation group:
the induced continuous automorphism of $S^1$ has degree
$+1$. This would contradict nonabelianness. These
arguments prove (C.37.9), including the topology of its
image in the intrinsic homeomorphism group.

For a line $B$, define $R^B_{r,t}$ by the logarithmic formula
in this simply connected group. For a circle, use the
unique based real lift of $f_t|_B$ to the cover of its
rotation group and scale that lift by $r$.
Its endpoint is zero by the exact winding comparison C.29.
On a point orbit define $R^B_{r,t}=1$.
These definitions are independent of the conjugacies in
(C.37.9), since Lie homomorphisms commute with exponential
maps and with based universal-cover lifts.
They agree with C.36 on every boundary orbit to which
that construction already applies.

\emph{Limits from the dense open union.}
Take
\[
 x_i\in O_i\subset U_2,\quad x_i\longrightarrow x\in X_2,
 \qquad (r_i,t_i)\longrightarrow(r,t).
                                                        \tag{C.37.10}
\]
Choose word isotopies $W_i$ from (C.37.7) representing
$R^{O_i}_{r_i,t_i}$ with their specified based lifts.
After a subsequence all word parameters converge.
Then $W_i$ and their inverses converge uniformly on compact
sets, uniformly in $u$, to a word isotopy $W$ from the
same family. Write $w=W(1)$ and $g=f_t$.
If $x$ is in an open orbit of $X_2$, eventually $x_i$ is
in that same open orbit, where (C.37.6) is continuous.
If $x$ is a point orbit, $w(x)=x$ for every possible
limiting word, which gives the required limit immediately.
It remains to identify $w$ on the intrinsic line or circle
$B=Gx$.

We will repeatedly use the following exact passage of
identities. For every fixed $h\in G$,
$hx_i\in O_i$ and $hx_i\to hx$. An identity of word
homeomorphisms on $O_i$, whose word isotopies converge,
therefore passes to every point $hx$ of $B$.
In particular
\[
                        [w|_B,g|_B]=1,                 \tag{C.37.11}
\]
because a fractional power commutes with its endpoint
on $O_i$.

Fix a positive integer $n$ and set $m_i=\lfloor nr_i\rfloor$.
In the universal covering group of $N_{O_i}$ the exact
residual is
\[
 \widetilde w_i^{\,n}\,
       \widetilde n_{O_i}(t_i)^{-m_i}
    =\exp\bigl((nr_i-m_i)
                  \log\widetilde n_{O_i}(t_i)\bigr).
                                                        \tag{C.37.12}
\]
Its exponent fraction lies in $[0,1)$.
By (C.37.2), the right side is again represented by one of
the compact family (C.37.7), with equality of based lifts.
Passing to a further subsequence for this fixed $n$,
the integer $m_i$ is constant, say $m_n$, and the residual
word isotopies converge to $Z_n$. Hence
\[
 (w|_B)^n(g|_B)^{-m_n}=z_n,\qquad
 |m_n-nr|\leq1,\qquad z_n=Z_n(1)|_B.                    \tag{C.37.13}
\]
All $z_n$ belong to the same compact set of restricted
three-word endpoints, independently of $n$.
No simultaneous subsequence for all $n$ is required.

For a line, (C.37.9) and (C.37.11) allow logarithms in
$\R$ or $\operatorname{Aff}^+(\R)$.
In the latter group the centralizer of every nonidentity
element is its exponential one-parameter subgroup, as is
seen directly from (C.37.4). Thus logarithms add for
commuting elements in both groups. The compact set in
(C.37.13) has bounded logarithms, giving
\[
 n\log(w|_B)-m_n\log(g|_B)=\log z_n,\qquad
 \sup_n\|\log z_n\|<\infty.
                                                        \tag{C.37.14}
\]
Divide by $n$ and let $n$ tend to infinity.
It follows that $\log(w|_B)=r\log(g|_B)$.
This includes $g|_B=1$, in which case boundedness forces
$w|_B=1$. The limiting word therefore has exactly the
required fractional action on the line.

\emph{The lifted residual bound on a circle.}
Identify the action in (C.37.9) on a circle with rotations,
and let $\theta(u)$ be the continuous real angle of
$f_u|_B$, normalized by $\theta(0)=0$.
The normalization uses rotations $\R/\Z$.
C.29 gives $\theta(1)=0$.
Every isotopy in (C.37.7), restricted to $B$, has a
canonical real endpoint angle
\[
 \epsilon\theta(t)+\theta(b)-\theta(a).
                                                        \tag{C.37.15}
\]
Let $A$ be this angle for $W$, let $T=\theta(t)$,
and let $C_n$ be the angle for $Z_n$.
Writing $M=\max_{0\leq u\leq1}|\theta(u)|$ gives
$|C_n|\leq3M$.
We must prove the real, rather than merely circle-valued,
identity
\[
                     nA-m_nT=C_n.                      \tag{C.37.16}
\]

The word isotopy
\[
 \Lambda_i(u)=W_i(u)^n f_{u t_i}^{-m_i}Z_i(u)^{-1}
                                                        \tag{C.37.17}
\]
restricts to a based loop in $N_{O_i}$.
Equation (C.37.12), with the specified lifts of $W_i$ and
$Z_i$, says that its universal-cover lift ends at the
identity. Hence that loop is null-homotopic in $N_{O_i}$.
Choose $h\in G$ with $hx\ne x$, and evaluate (C.37.17)
at the two points $x_i,hx_i$ of $O_i$.
For large $i$ they are distinct. The resulting loop in
$\operatorname{Conf}_2(\R^2)$ has difference-direction
winding zero. The word isotopies converge uniformly on
compact sets to
$\Lambda(u)=W(u)^n f_{ut}^{-m_n}Z_n(u)^{-1}$.
The two-point loops therefore converge uniformly to its
evaluation at $x,hx$. This limiting pair stays distinct,
so winding is preserved.
On the Jordan circle $B$ its winding, by C.29, is exactly
$nA-m_nT-C_n$ (up to a common orientation sign).
It is zero, proving (C.37.16).

Using (C.37.13) and $|C_n|\leq3M$ now gives
\[
        |A-rT|\leq\frac{3M+|T|}{n}\quad(n\geq1).
                                                        \tag{C.37.18}
\]
Consequently $A=rT$. This determines the based lifted
fractional action even when $f_t|_B=1$ but $\theta(t)$
is a nonzero integer. Thus no hidden winding survives
the limit from the open orbits.

\emph{Joint continuity, inverses, and equivariance.}
We have shown that every convergent subsequence of word
parameters in (C.37.10) gives the same endpoint
$R^B_{r,t}(x)$. Compactness of (C.37.7) proves
\[
          R^{O_i}_{r_i,t_i}(x_i)
                  \longrightarrow R^B_{r,t}(x).
                                                        \tag{C.37.19}
\]
This holds for every sequence from the dense set $U_2$
approaching every point of $X_2$, with parameters also
varying. It proves joint continuity on all of $X_2$.
Explicitly, for a sequence of points of $X_2$, approximate
each by a point of $U_2$ so closely that both its distance
and its image distance for that fixed parameter are
less than $1/i$. Such approximations exist by
(C.37.19) with parameters fixed. Applying (C.37.19) to
the chosen approximating sequence gives the full
continuity assertion.

The inverse loop $f_t^{-1}$ has the same group $G$,
the same sets $U_2,X_2$, and the same actual operator
images. Its construction gives precisely
$(R^B_{r,t})^{-1}$, by inversion of the based lifts.
Applying the preceding argument to it proves joint
continuity of the inverse family as well. Each member
is an orbitwise bijection, hence a homeomorphism of
$X_2$. The identities $R_{0,t}=1$, $R_{1,t}=f_t|_{X_2}$,
and $R_{r,0}=R_{r,1}=1$ give the based contraction.

Finally every $h\in H$ permutes the open $G$-orbits and
conjugates their actual closed operator images.
It therefore preserves $U_2$ and $X_2$.
C.36 gives exact transported equivariance on $U_2$.
The established joint continuity extends that equality
to its closure $X_2$. This proves all assertions.
\end{proof}

\paragraph{The remaining global boundary requirement.}
C.37 treats the whole closed set $X_2$, including infinite
accumulation of its open orbits. Its boundary contraction
is the based Lie-group fractional action forced by
(C.37.14) and (C.37.16). Extending this family to the
complement of $X_2$, compatibly with the dimension-zero
and dimension-one proper operator cases, remains a
separate requirement. Even where that complement
contains only nonopen orbits, its median contraction
from C.34 has not been identified with this boundary
family. C.35 proves that imposing the fixed median
prescription can violate the actual restriction map.
Thus the uniform estimate proved here closes the
dimension-two accumulation step, while the global
centralizer conclusion in C.12 remains open.
